% ===========================================================================
%  figure/tikz/pipeline-motore.tex
%  Che cosa accade a ogni barra del backtest, in ordine.
% ===========================================================================
\begin{tikzpicture}[node distance = 4mm]

  \node[archivio]                           (dati)  {parquet\\OHLCV};
  \node[motore, right = 11mm of dati, minimum width = 26mm]
    (ind)  {indicatori\\{\scriptsize pre-calcolati}};
  \node[motore, right = 11mm of ind, minimum width = 26mm]
    (strat) {strategie\\{\scriptsize \codT{generate\_signals}}};
  \node[motore, right = 11mm of strat, minimum width = 26mm]
    (allo) {allocator\\{\scriptsize dai segnali ai pesi}};

  \node[motore, below = 9mm of allo, minimum width = 26mm]
    (risk) {rischio\\{\scriptsize leva, vol target}};
  \node[motore, left = 11mm of risk, minimum width = 26mm]
    (exec) {esecuzione\\{\scriptsize slippage, costi}};
  \node[motore, left = 11mm of exec, minimum width = 26mm]
    (port) {portafoglio\\{\scriptsize posizioni, equity}};
  \node[archivio, left = 11mm of port]      (out)   {parquet\\PnL};

  \draw[flusso] (dati)  -- (ind);
  \draw[flusso] (ind)   -- (strat);
  \draw[flusso] (strat) -- (allo) node[etichetta, above, pos = 0.5] {\codT{Signal}};
  \draw[flusso] (allo)  -- (risk);
  \draw[flusso] (risk)  -- (exec);
  \draw[flusso] (exec)  -- (port) node[etichetta, above, pos = 0.5] {\codT{Fill}};
  \draw[flusso] (port)  -- (out);

  \node[passo] at ([shift = {(-3mm,3mm)}]ind.north west)   {1};
  \node[passo] at ([shift = {(-3mm,3mm)}]strat.north west)  {2};
  \node[passo] at ([shift = {(-3mm,3mm)}]allo.north west)   {3};
  \node[passo] at ([shift = {(3mm,3mm)}]risk.north east)    {4};
  \node[passo] at ([shift = {(3mm,3mm)}]exec.north east)    {5};
  \node[passo] at ([shift = {(3mm,3mm)}]port.north east)    {6};

  \node[nota, below = 6mm of port, text width = 92mm, align = center]
    {I passi 1--3 decidono con i dati fino alla chiusura di $T-1$.
     I passi 4--6 eseguono all'apertura di $T$.};

\end{tikzpicture}
